% Generated from validated Article Draft Result. Do not hand-edit; revise the structured draft instead.
\section{生成から行動へ――RAG、Tool Use、Computer Useの接続}
\noindent\textbf{2024年上期には、モデルが知識を返すだけでなく、検索・外部tool・computer environmentへ接続して仕事を進めるための設計が前面に出た。Command R、SWE-agent、OSWorldを軸に、RAGからactionまでの段階差を整理する。}\autocite{src-1c9f1d22c176,src-70b5d4b3bedd,src-77e35e5f22be}

Command R familyが示したのは、長文脈だけではなくretrievalやtool useを前提としたmodel/application設計である。RAGは外部情報を取り込む層、tool useは外部機能を呼び出す層、agent的実行は複数stepで環境を変化させる層であり、同じ「外部接続」でも失敗モードと制御条件は異なる。\autocite{src-77e35e5f22be}


SWE-agentはsoftware engineering taskをenvironment interactionとして扱い、modelがrepositoryやtoolsへ働きかける評価系を提示した。ここではcode generation能力だけでなく、状態を観察し、toolを選び、失敗から修正するloop全体が性能を左右する。単発の生成品質からexecution trajectoryへ評価単位が広がった点が重要である。\autocite{src-1c9f1d22c176}


OSWorldはcomputer useを現実的なOS環境に近いbenchmarkとして扱い、GUI操作を含むaction能力を独立した研究対象にした。Quiet-STaRのような内部推論研究はaction systemそのものではないが、複数stepの判断を支えるreasoning layerの探索として接続する。Model Specは、こうした行動可能なmodelをどのbehavioral rulesで制御するかという別レイヤの課題を浮かび上がらせる。\autocite{src-12fa04e64d24,src-70b5d4b3bedd,src-d1406200c99e}


\begin{claimboundary}
benchmark上でtoolやGUIを操作できることと、実運用で安全かつ安定して権限を委譲できることは同義ではない。外部actionには誤操作、権限境界、prompt injection、状態認識、rollbackなどmodel単体にはないfailure surfaceが加わる。本稿ではbenchmark resultをdeployed agentの信頼性保証へ読み替えない。\autocite{src-1c9f1d22c176,src-2f7e28810d05,src-70b5d4b3bedd,src-77e35e5f22be,src-edd253654067}
\end{claimboundary}

2024-H1の変化は、「よりよい文章を生成するモデル」から「情報を取り込み、機能を呼び、環境へ作用するsystem」への入口が見えたことにある。RAG、tool call、computer useは同一機能ではないが、能力の単位をmodel outputからclosed-loop executionへ拡張するという点で同じ方向を向いていた。\autocite{src-1c9f1d22c176,src-70b5d4b3bedd,src-77e35e5f22be,src-d1406200c99e}
